% !TeX root = sections_1_3_UK_simple.tex
% Проста українська редакція розділів 1--3 для порівняння до/після.
\begingroup\color{black}
\hypersetup{citecolor=black,linkcolor=black}

\section{Вступ}
\label{sec:introduction}

Інфекція поширюється як усередині спільнот, так і між ними. Усередині спільноти важливо, хто саме з ким контактує \citep{keeling2005,pastorsatorras2015,kiss2017}. Про зв'язки між спільнотами часто відомо менше: наприклад, лише кількість пасажирів або поїздок на роботу \citep{balcan2009,barbosa2018human}. Ми досліджуємо модель, яка зберігає окремі контакти всередині спільнот, а зв'язки між ними описує загальними вагами.

Контактні мережі можна будувати за даними про населення, його переміщення та контакти \citep{eubank2004,danon2011}. Зв'язки між спільнотами можна оцінювати за транспортними й міграційними даними або даними мобільних операторів \citep{barbosa2018human,chang2021}. Такі дані показують, наскільки сильно пов'язані дві спільноти, але не завжди вказують конкретних людей, через яких може перейти інфекція. Вузли, що мають контакти з іншою спільнотою, далі називаємо вузлами-мостами.

Якщо відомі всі контакти, можна моделювати випадкові зараження методом Gillespie. Він точно відтворює заданий неперервний марковський ланцюг (CTMC) \citep{gillespie1977exact,gibson2000}. Існують також ефективні алгоритми для мереж зі змінними в часі контактами та для мережевого моделювання епідемій \citep{vestergaard2015,kuryliak2025simulating}. Їм потрібні конкретні контакти, через які передається інфекція. Метапопуляційні моделі натомість описують спільноти кількостями людей у різних станах і зв'язують їх загальними потоками \citep{colizza2007reaction,colizza2008,balcan2010modeling,ajelli2010}. Ми поєднуємо ці підходи: моделюємо окремі вузли всередині спільнот і використовуємо загальні ваги для передавання між ними.

Ідею двомасштабного симулятора представлено в попередній конференційній роботі \citep{kuryliak2025towards}. Кожна спільнота має свій локальний симулятор. Спільний планувальник, який називаємо макрорівнем, визначає зараження між спільнотами за їхніми поточними станами та вагами зв'язків. Зараження, що надходить з іншої спільноти, називаємо імпортом інфекції. У цій статті описуємо метод і досліджуємо його похибку: що змінюється, коли замість справжніх міжспільнотних контактів відома лише їхня загальна вага, а імпорт потрапляє на випадково вибраний сприйнятливий вузол?

Головне питання: коли цього спрощення достатньо, щоб правильно оцінити час поширення інфекції? Ми порівнюємо часи її першого прибуття до спільнот і розподіли цих часів за багатьма запусками. Також перевіряємо, чи могло зараження відбутися через справжні контакти. Близькі часи ще не гарантують можливого шляху передавання: двомасштабна модель може заразити іншу спільноту, поки всі справжні вузли-мости в спільноті-джерелі ще сприйнятливі. Для перевірки ми використовуємо синтетичні мережі, у яких знаємо всі контакти. Повний Gillespie отримує всю мережу, а двомасштабний метод --- лише локальні мережі та загальні ваги міжспільнотних зв'язків.

Стаття має чотири основні результати.
\begin{enumerate}
\item Описуємо, як працює симулятор: кожна спільнота моделюється окремо, а спільний планувальник визначає зараження між ними. Пояснюємо, як він узгоджує локальні симулятори на часових інтервалах фіксованої довжини (Розділ~\ref{sec:method}, Алгоритм~\ref{alg:scheduler}).
\item Показуємо, яку мережу метод моделює точно, якщо визначати час кожної події без наближення часовою сіткою. Це \emph{сурогатна мережа}: справжні контакти між двома спільнотами замінено слабкими зв'язками між усіма парами їхніх вузлів зі збереженням загальної ваги (Твердження~\ref{prop:surrogate}). Пояснюємо джерела похибки цієї заміни. Для одного передавання назовні доводимо, що за випадкового рівномірного розміщення мостів такий метод схильний передавати інфекцію раніше за повний Gillespie; точний зміст цього порівняння наведено в Твердженні~\ref{prop:ordering}.
\item Порівнюємо метод із повним Gillespie для одного передавання інфекції назовні та для ланцюгів спільнот. У наших тестах за швидкого поширення всередині спільнот розподіли часу прибуття зазвичай близькі. За повільнішого внутрішнього поширення двомасштабний метод систематично дає раніше прибуття. Випадкове переставляння міток методів у результатах запусків допомагає оцінити, чи можна пояснити відмінності скінченною кількістю запусків (Розділи~\ref{sec:design} і~\ref{sec:results}).
\item Пропонуємо показник, який допомагає пояснити точність у наших тестах. Він порівнює розкид часів зараження всередині спільноти із середнім часом передавання через одне міжспільнотне ребро. За малих значень розподіли часу прибуття зазвичай узгоджуються; за більших метод дає раніше прибуття. Це спостереження для досліджених мереж, а не загальний закон похибки (Розділ~\ref{sec:res_descriptor}).
\end{enumerate}

У всіх експериментах вузол після зараження залишається інфікованим: це модель «сприйнятливий--інфікований» (SI). Мережі статичні й синтетичні, а крок синхронізації фіксований. Далі описуємо попередні праці (Розділ~\ref{sec:related}), метод (Розділ~\ref{sec:method}) та інформацію, яку він зберігає і втрачає (Розділ~\ref{sec:analysis}). Експерименти й результати наведено в Розділах~\ref{sec:design} і~\ref{sec:results}, а застосування та обмеження --- у Розділі~\ref{sec:discussion}.

\clearpage
\section{Пов'язані праці}
\label{sec:related}

\paragraph{Моделювання випадкових подій.}
Прямий метод Gillespie визначає час наступної події та вибирає її за поточними інтенсивностями --- величинами, що задають частоту подій \citep{gillespie1977exact}. Метод наступної реакції веде окремий годинник для кожної можливої події \citep{gibson2000}. Обидва точно моделюють заданий CTMC. Якщо інтенсивність змінюється, час події можна визначати, накопичуючи її до випадкового порога \citep{ethier1986}. \citet{anderson2007} застосував цю ідею до систем зі змінними інтенсивностями та затримками. Так само працює наш годинник міжспільнотних заражень. Для ефективного мережевого моделювання використовують, зокрема, часовий алгоритм Gillespie та оновлення лише тих інтенсивностей, які змінила остання подія \citep{vestergaard2015,kuryliak2025simulating}. Наші локальні симулятори використовують реалізацію з \citet{kuryliak2025simulating}.

Багатомасштабні методи для хімічних реакцій можуть наближено описувати швидкі реакції, щоб прискорити моделювання \citep{haseltine2002,cao2005slow}. Наше спрощення стосується контактів: локальні зараження моделюються окремо, але конкретні міжспільнотні ребра невідомі. Співвідношення часу поширення всередині спільноти й часу передавання між спільнотами допомагає пояснити наслідки цього спрощення.

\paragraph{Моделі з кількома спільнотами.}
Метапопуляційні моделі рахують людей у різних станах у кожній спільноті й описують переміщення між спільнотами. Так досліджують глобальне поширення та умови проникнення інфекції \citep{colizza2007reaction,colizza2008,balcan2009,balcan2010modeling}. Дані про переміщення також поєднують із локальними моделями передавання \citep{chang2021}. Агентні моделі описують окремих людей у всій системі \citep{eubank2004,kerr2021}. \citet{ajelli2010} порівняли ці підходи. Наш метод зберігає окремі вузли та контакти всередині спільнот, а зараження між ними визначає за загальними вагами зв'язків і поточними кількостями вузлів у різних станах.

\paragraph{Два рівні перемішування та спільне моделювання.}
Деякі епідемічні моделі враховують окремі контакти в малих групах, а на ширшому рівні припускають рівномірне перемішування \citep{ball1997}. Подібне припущення має наша сурогатна мережа: між двома зв'язаними спільнотами загальну вагу рівномірно розподілено між усіма парами вузлів.

Метод також пов'язаний зі спільним моделюванням (co-simulation): окремі симулятори працюють за узгодженими часовими інтервалами й обмінюються даними \citep{gomes2018cosim}. У нас кожен локальний симулятор повідомляє кількості інфікованих і сприйнятливих вузлів. Планувальник використовує їх для визначення часу міжспільнотного зараження та вибору його цілі. Йому не потрібно знати всі внутрішні кроки локальних симуляторів.

\clearpage
\section{Метод двомасштабного моделювання}
\label{sec:method}

Усі експерименти використовують одну реалізацію двомасштабного симулятора. Вона узгоджує локальні симулятори на часових інтервалах фіксованої довжини та наближено визначає час зараження між спільнотами. Для пояснення похибки також розглядаємо теоретичну версію з точним часом кожної події. Далі називаємо її \emph{ідеальним процесом}. Це двомасштабний метод без чисельної похибки часової сітки. Він усе ще використовує усереднені зв'язки між спільнотами, тому точно моделює сурогатну мережу, а загалом не початкову мережу зі справжніми контактами. Нижче описуємо обидві мережі та роботу симулятора.

\subsection{Повномережева еталонна модель}
\label{sec:reference}

Нехай $G=(V,E)$ --- контактна мережа. Її вузли поділено на $M$ спільнот, які не перетинаються:
\begin{equation}
\label{eq:partition}
V=V_1\cup V_2\cup\cdots\cup V_M,
\qquad N_i=|V_i|.
\end{equation}
Тут $V_i$ --- множина вузлів спільноти $i$, а $N_i$ --- кількість її вузлів. Позначаємо через $E_{ii}$ контакти всередині спільноти $i$, а через $E_{ij}$ --- контакти від вузлів спільноти $i$ до вузлів спільноти $j$. Контакт представляємо як напрямлений канал передавання. Тому ненапрямлений контакт дає два канали --- по одному в кожному напрямку.

Кожен вузол є сприйнятливим ($S$), інфікованим ($I$) або видаленим ($R$). Повний стан мережі в час $t$ позначаємо $X(t)$, а стан спільноти $i$ --- $X_i(t)$. Множини сприйнятливих та інфікованих вузлів цієї спільноти позначаємо $S_i(t)$ та $I_i(t)$. Їхні кількості --- $s_i(t)$ та $y_i(t)$, а частки --- $\sbar_i(t)=s_i(t)/N_i$ та $\ibar_i(t)=y_i(t)/N_i$. Далі $x$ означає конкретний стан мережі; $x^{v:I}$ --- той самий стан після зараження вузла $v$, а $x^{u:R}$ --- після видалення вузла $u$.

Повна модель має два прості правила. Сприйнятливий вузол $v$ отримує інфекцію від кожного інфікованого сусіда $u$ з інтенсивністю $\beta_{uv}$. Інфікований вузол видаляється з інтенсивністю $\gamma$. Для вузлів у відповідних станах:
\begin{equation}
\label{eq:full_generator}
q_v(x)=\sum_{\substack{u:(u,v)\in E\\x_u=I}}\beta_{uv}
\quad (x_v=S),
\qquad r_u(x)=\gamma\quad (x_u=I).
\end{equation}
Тут $q_v$ --- сумарна інтенсивність зараження вузла $v$, а $r_u$ --- інтенсивність видалення вузла $u$. Для несприйнятливого вузла $q_v=0$; для неінфікованого вузла $r_u=0$. Інтенсивність визначає частоту події: наприклад, за сталої інтенсивності $q$ середній час очікування дорівнює $1/q$. Ці правила також визначають генератор CTMC $\mathcal{L}_G$, використаний далі; ця назва не додає нових правил епідемії. Для міжспільнотного каналу записуємо $\beta_{uv}=\beta_{\mathrm{out}}w_{uv}$, де $w_{uv}$ --- вага контакту. Точне моделювання цих повномережевих інтенсивностей позначаємо \Full{}. Усі експерименти в статті мають динаміку SI, тобто $\gamma=0$.

\subsection{Агрегований інтерфейс}
\label{sec:interface}

Двомасштабний симулятор \TS{} отримує всі локальні графи $G_i=(V_i,E_{ii})$ та їхні інтенсивності. Конкретні контакти між спільнотами йому невідомі. Замість них для кожного напрямку $i\to j$ задано загальну вагу:
\begin{equation}
\label{eq:weight}
W_{ij}=\sum_{(u,v)\in E_{ij}}w_{uv},\qquad i\neq j.
\end{equation}
Якщо всі ваги дорівнюють одиниці, $W_{ij}$ --- кількість каналів від $V_i$ до $V_j$. У прикладній задачі цю величину можна оцінити за транспортними, міграційними даними, даними мобільності або контактів. Макрорівень також отримує поточні кількості сприйнятливих та інфікованих вузлів від кожного локального симулятора. Він не знає, які вузли є кінцями справжніх міжспільнотних ребер. Основні відмінності наведено в Таблиці~\ref{tab:information}; Рисунок~\ref{fig:schematic} показує їх на схемі.

\begin{table}[tbp]
\color{black}\centering
\caption{\textcolor{black}{Дані та правила двох моделей. В усіх експериментах \TS{} вибирає цільовий вузол рівномірно серед сприйнятливих вузлів.}}
\label{tab:information}
\small
\begin{tabular}{>{\raggedright\arraybackslash}p{0.42\linewidth}>{\raggedright\arraybackslash}p{0.23\linewidth}>{\raggedright\arraybackslash}p{0.25\linewidth}}
\toprule
Дані або дія & \Full{} & \TS{}\\
\midrule
Локальні контакти, інтенсивності та стани вузлів & відомі & відомі локальним симуляторам\\
Конкретні міжспільнотні ребра та вузли-мости & відомі & невідомі\\
Загальна міжспільнотна вага $W_{ij}$ & отримана з ребер & задана безпосередньо\\
Інтенсивність міжспільнотного зараження & зі станів кінців справжніх ребер & з $\ibar_i$, $\sbar_j$ та $W_{ij}$\\
Цільовий вузол зараження & кінець ребра, яке передало інфекцію & сприйнятливий вузол, вибраний за правилом\\
Кількість можливих SI-заражень через набір каналів & не більше кількості каналів & не більше $s_j$ цільових заражень\\
\bottomrule
\end{tabular}
\end{table}

\begin{figure}[tbp]
\color{black}\centering
\begin{tikzpicture}[
  x=1cm,y=1cm,
  font=\scriptsize,
  >=Latex,
  person/.style={circle,draw=black!60,fill=blue!18,minimum size=4.2mm,inner sep=0pt},
  infected/.style={person,fill=orange!70},
  bridge/.style={person,very thick,draw=black},
  bridgeinf/.style={infected,very thick,draw=black},
  commbox/.style={draw=black!45,dashed,rounded corners=2pt,inner sep=3pt},
  macrobox/.style={draw=black!70,rounded corners=2pt,fill=black!5,align=center,inner sep=3pt},
  link/.style={draw=black!55},
  weak/.style={draw=black!24,line width=.25pt},
  flow/.style={-{Latex[length=1.7mm]},draw=black!65}
]
% (a) Full microscopic interface
\node[font=\bfseries] at (0,2.25) {(a) Full network};
\node[infected] (a1) at (-1.65,1.20) {};
\node[person]   (a2) at (-1.55,0.35) {};
\node[bridgeinf](a3) at (-0.75,0.82) {};
\node[person]   (a4) at (-0.75,1.55) {};
\node[bridge]   (b1) at (0.65,1.15) {};
\node[person]   (b2) at (1.45,1.55) {};
\node[person]   (b3) at (1.55,0.55) {};
\node[bridge]   (b4) at (0.70,0.30) {};
\draw[link] (a1)--(a2)--(a3)--(a1) (a1)--(a4)--(a3);
\draw[link] (b1)--(b2)--(b3)--(b4)--(b1) (b1)--(b3);
\draw[very thick] (a3)--(b1);
\draw[very thick] (a3)--(b4);
\node[commbox,fit=(a1)(a2)(a3)(a4)] {};
\node[commbox,fit=(b1)(b2)(b3)(b4)] {};
\node at (-1.2,-0.10) {$V_i$};
\node at (1.15,-0.10) {$V_j$};

% (b) Two-scale representation
\node[font=\bfseries] at (5.0,2.25) {(b) Two-scale simulator};
\node[macrobox] (macro) at (5.0,1.55) {macro interface\\$W_{ij}$, $(\ibar_i,\sbar_j)$\\$\lambda_{ij}$ + placement};
\node[macrobox,minimum width=2.15cm,minimum height=.70cm] (li) at (3.65,0.35) {local CTMC\\$G_i$};
\node[macrobox,minimum width=2.15cm,minimum height=.70cm] (lj) at (6.35,0.35) {local CTMC\\$G_j$};
\draw[flow] (li.north) -- node[left] {summaries} (macro.south west);
\draw[flow] (macro.south) -- node[right] {import} (lj.north);
\draw[flow] (lj.east) -- ++(0,1.0) -| node[pos=1.08,right] {summaries} (macro.east);
\node at (5.0,-0.35) {node-level dynamics retained locally};

% (c) Surrogate interface
\node[font=\bfseries] at (10.25,2.25) {(c) Surrogate network};
\node[infected] (c1) at (8.85,1.45) {};
\node[person]   (c2) at (8.75,0.55) {};
\node[person]   (c3) at (9.55,1.05) {};
\node[person]   (c4) at (9.45,0.20) {};
\node[person]   (d1) at (11.00,1.45) {};
\node[person]   (d2) at (10.90,0.55) {};
\node[person]   (d3) at (11.70,1.05) {};
\node[person]   (d4) at (11.60,0.20) {};
\foreach \u in {c1,c2,c3,c4}{\foreach \v in {d1,d2,d3,d4}{\draw[weak] (\u)--(\v);}}
\draw[link] (c1)--(c2)--(c4)--(c3)--(c1);
\draw[link] (d1)--(d2)--(d4)--(d3)--(d1);
\node[commbox,fit=(c1)(c2)(c3)(c4)] {};
\node[commbox,fit=(d1)(d2)(d3)(d4)] {};
\node[align=center] at (10.25,-0.35) {$W_{ij}/(N_iN_j)$ per\\cross-community channel};
\end{tikzpicture}
\caption{\textcolor{black}{Три подання зв'язку між спільнотами. (a) У повній мережі інфекція переходить лише через показані фізичні ребра та їхні вузли-мости. (b) Двомасштабний метод зберігає окремий локальний симулятор для кожної спільноти. Макрорівень отримує кількості вузлів у різних станах, використовує $W_{ij}$ для генерування міжспільнотного зараження та вибирає його ціль рівномірно серед сприйнятливих вузлів. (c) Ідеальний двомасштабний процес еквівалентний процесу на сурогатній мережі, де кожна напрямлена пара вузлів між двома спільнотами має слабкий зв'язок. Загальна вага всіх цих зв'язків дорівнює $W_{ij}$. Помаранчеві вузли --- інфіковані, сині --- сприйнятливі.}}
\label{fig:schematic}
\end{figure}

\subsection{Локальні симулятори}
\label{sec:local}

Кожна спільнота має свій симулятор Gillespie для графа $G_i$. Він моделює зараження та видалення за~\eqref{eq:full_generator}, враховуючи лише внутрішні контакти $E_{ii}$. Після події оновлюються лише ті інтенсивності, які вона змінила, як у \citet{kuryliak2025simulating}. Кожен симулятор має власний потік випадкових чисел.

Планувальник задає час, до якого локальний симулятор має виконати свої події. У цей момент він може додати імпорт: заразити вибраний сприйнятливий вузол. Після цього локальні інтенсивності змінюються. Симулятор зберігає невикористану частину випадкового порога наступної локальної події й перераховує час очікування за новою сумарною інтенсивністю. Подію та її вузол вибирає за оновленими інтенсивностями. Новий поріг генерує лише тоді, коли збереженого немає. Локальні пороги й поріг міжспільнотної події окремі; останній описано в Розділі~\ref{sec:timing}. За заданих часів і цілей імпорту локальний симулятор точно виконує правила своєї мережі. Похибка зв'язування спільнот залежить від того, як планувальник визначає імпорти та вибирає їхні цілі.

\subsection{Макрорівень}
\label{sec:macro}

Планувальник використовує частки вузлів у різних станах і вагу $W_{ij}$. Інтенсивність передавання зі спільноти $i$ до спільноти $j$ дорівнює
\begin{equation}
\label{eq:macro_rate}
\lambda_{ij}(t)=\beta_{\mathrm{out}}W_{ij}\ibar_i(t)\sbar_j(t),
\qquad \Theta(t)=\sum_{i\neq j}\lambda_{ij}(t).
\end{equation}
Тут $\beta_{\mathrm{out}}W_{ij}$ задає загальну силу зв'язку, $\ibar_i$ --- частку інфікованих у джерелі, а $\sbar_j$ --- частку сприйнятливих в одержувачі. Що більше цих вузлів, то вища інтенсивність передавання. Сума $\Theta$ враховує всі напрямки між спільнотами.

Коли така подія відбувається, спочатку вибирають спільноту-джерело та спільноту-одержувача $(i,j)$ з імовірністю $\lambda_{ij}/\Theta$. Потім у спільноті $j$ рівномірно вибирають один сприйнятливий вузол і заражають його. Якщо таких вузлів 20, кожен має імовірність $1/20$. Вибір не враховує справжні вузли-мости, оскільки агрегований інтерфейс їх не знає. У подальших формулах цю імовірність позначаємо $\Pi_{ij}(v\mid x)=1/s_j$ для $v\in S_j$, де $x$ --- поточний стан мережі. Отже, $\Pi$ просто позначає імовірність вибору конкретного сприйнятливого вузла. Інші правила можливі, якщо є інформація про вузли-мости. Додатковий калібрувальний множник у реалізації дорівнює одиниці в усіх експериментах.

Формула усереднює передавання між спільнотами. За заданих кількостей інфікованих і сприйнятливих вона дорівнює очікуваній інтенсивності, якщо кінці міжспільнотних ребер вибрати рівномірно й випадково. Для конкретної мережі інтенсивність може бути іншою: важливі стани саме вузлів із міжспільнотними контактами та історія заражень. Для сурогатної мережі зі зв'язками «усі з усіма» формула точна (Розділ~\ref{sec:analysis}). Значення $\Theta$ оновлюється після подій, що змінюють кількості вузлів у різних станах.

\subsection{Час подій та синхронізація}
\label{sec:timing}

Після міжспільнотного зараження в час $t_0$ генеруємо випадковий поріг $E\sim\Exp(1)$ --- експоненційне число із середнім один. Від цього моменту накопичуємо загальну інтенсивність $\Theta$. Наступне міжспільнотне зараження відбувається в перший момент $t^\ast$, коли накопичена величина досягає порога:
\begin{equation}
\label{eq:hazard_integral}
\int_{t_0}^{t^\ast}\Theta(s)\,ds=E.
\end{equation}
Інтеграл --- це площа під кривою інтенсивності. Вона зростає до порога $E$; якщо його не досягнуто, подія не відбувається. Так визначають час події за змінної інтенсивності \citep{ethier1986,anderson2007}. За сталої $\Theta$ час очікування дорівнює $E/\Theta$. Якщо взяти $E=-\ln U$, де $U$ рівномірно розподілене на $(0,1)$, отримуємо звичну формулу Gillespie. Коли локальна подія змінює $\Theta$, зберігаємо поріг і вже накопичену величину, а далі додаємо внесок нової інтенсивності. Новий поріг генеруємо після міжспільнотного зараження. Цей принцип використано і в конференційній роботі. Ідеальний процес визначає момент досягнення порога точно; наша реалізація оцінює його за часовою сіткою. Ідеальний процес точний для сурогатної мережі (Розділ~\ref{sec:surrogate}).

У реалізації час поділено на інтервали довжини $\tau$. Час міжспільнотного зараження може лежати всередині інтервалу: його не округлюють до точки сітки. Для оцінки використовуємо інтенсивність у середині інтервалу, як у конференційному описі. Позначимо через $A$ уже накопичену величину. На поточному інтервалі $[t,t_g)$ копіюємо локальні симулятори разом зі станами їхніх генераторів випадкових чисел. Копії виконуємо до середини інтервалу й за їхніми станами обчислюємо $\Theta_{\mathrm{mid}}$. Оцінка приросту дорівнює $\Theta_{\mathrm{mid}}(t_g-t)$. Якщо $A$ разом із цим приростом ще менше за $E$, справжні симулятори виконуються до $t_g$. Інакше оцінюємо час зараження:
\begin{equation}
\label{eq:interpolation}
t^\ast=t+\frac{E-A}{\Theta_{\mathrm{mid}}}.
\end{equation}
Залишок до порога дорівнює $E-A$. Ділення на оцінену інтенсивність дає час, потрібний для його накопичення. Наприклад, якщо $E=1$, $A=0.4$ і $\Theta_{\mathrm{mid}}=0.2$, час очікування дорівнює трьом одиницям часу. Планувальник виконує справжні симулятори до $t^\ast$ і за їхніми станами вибирає напрямок передавання та сприйнятливий вузол-одержувач. Потім заражає цей вузол. Повну послідовність наведено в Алгоритмі~\ref{alg:scheduler}.

\begin{algorithm}[H]
\color{black}
\caption{\textcolor{black}{Один запуск двомасштабного симулятора з узгодженням за часовою сіткою.}}
\label{alg:scheduler}
\small
\fbox{\begin{minipage}{0.94\linewidth}
\textbf{Вхід:} локальні графи $G_i$, ваги $W_{ij}$, інтенсивності, крок $\tau$, горизонт $T_{\mathrm{end}}$ та початковий стан.
\begin{enumerate}
\item Установити $t=0$ і $A=0$; згенерувати $E\sim\Exp(1)$.
\item Поки $t<T_{\mathrm{end}}$ й інфекція може поширюватися, встановити $t_g=\min(t+\tau,T_{\mathrm{end}})$ та обробити інтервал до $t_g$.
\item Скопіювати кожен локальний симулятор разом зі станом генератора випадкових чисел. Просунути копії до середини залишку інтервалу та обчислити $\Theta_{\mathrm{mid}}$ за~\eqref{eq:macro_rate}.
\item Якщо $A+\Theta_{\mathrm{mid}}(t_g-t)<E$, просунути справжні симулятори до $t_g$, додати цю величину до $A$, встановити $t=t_g$ і повернутися до кроку 2.
\item Інакше обчислити $t^\ast$ за~\eqref{eq:interpolation} та просунути справжні симулятори до $t^\ast$.
\item Вибрати $(i,j)$ з імовірністю $\lambda_{ij}(t^\ast)/\Theta(t^\ast)$. Рівномірно вибрати $v$ з $S_j(t^\ast)$ та імпортувати зараження у $v$. Оновити локальний симулятор одержувача, зберігаючи невикористаний експоненційний поріг і вибираючи локальну подію за новими інтенсивностями.
\item Установити $t=t^\ast$ і $A=0$, згенерувати нове $E\sim\Exp(1)$ та так само обробити залишок поточного інтервалу.
\end{enumerate}
\end{minipage}}
\end{algorithm}

Копії використовують той самий стан генератора випадкових чисел, що й справжні симулятори. Тому до першого імпорту вони відтворюють ту саму локальну траєкторію. Якщо на інтервалі немає локальних подій, інтенсивність стала й оцінка за серединою інтервалу точна. Якщо такі події є, оцінка може використати стан, пізніший за оцінений час імпорту, не більш ніж на $\tau/2$. Менший крок зменшує цей часовий зсув; у границі $\tau\to0$ похибка сітки зникає.

Щоб обійтися без часової сітки, локальні симулятори мають повідомляти часи своїх наступних подій. Тоді планувальник може обробляти події в часовому порядку й точно накопичувати $\Theta$, яка стала між подіями. Імпорт застосовується в момент його виникнення, після чого локальні інтенсивності оновлюються. Такий варіант усуває похибку часової синхронізації, але зберігає усереднений опис міжспільнотних зв'язків. В експериментах використано варіант із часовою сіткою та кроком $\tau=1$.

\subsection{Обчислювальна структура та застосування}
\label{sec:cost}

\Full{} зберігає всі контакти, зокрема між спільнотами. Обсяг роботи після події залежить від кількості сусідів відповідного вузла \citep{kuryliak2025simulating}. \TS{} зберігає локальні графи та матрицю міжспільнотних ваг розміру $M\times M$. Таке скорочення даних обговорено в конференційній роботі \citep{kuryliak2025towards}. Якщо міжспільнотних ребер багато, заміна їх загальними вагами може зменшити обсяг даних. Якщо ребер мало, економія може бути незначною. Локальні графи в будь-якому разі потрібно зберігати.

Локальні симулятори можна виконувати паралельно до спільного часу $t_g$ або $t^\ast$, який уже визначив планувальник. Копії для оцінки інтенсивності також можна виконувати паралельно. Потім потрібно зібрати стани спільнот і застосувати імпорт, якщо він має відбутися. У нашій реалізації локальні симулятори виконуються послідовно; можливе паралельне виконання ще не дає підстав заявляти про прискорення. Для різних спільнот також можна задавати різні локальні налаштування.

Копіювання станів і синхронізація потребують додаткової роботи, тому метод не гарантує меншого використання пам'яті або швидшого виконання. Його можна застосовувати, коли конкретні міжспільнотні контакти невідомі або їх надто багато для явного зберігання. Час виконання наведено в Розділі~\ref{sec:disc_cost}. Основне питання цієї статті --- точність такого спрощення.

\clearpage
\endgroup

